\documentclass[11pt]{article}
\usepackage[margin=1in]{geometry}
\usepackage{amsmath,amssymb,amsthm,mathtools}
\IfFileExists{lmodern.sty}{\usepackage{lmodern}}{}
\usepackage[expansion=false]{microtype}
\usepackage{booktabs}
\usepackage{enumitem}
\usepackage{xcolor}
\usepackage[colorlinks=true,linkcolor=blue!50!black,citecolor=green!35!black,urlcolor=blue!50!black]{hyperref}
\usepackage[nameinlink]{cleveref}

\newtheorem{theorem}{Theorem}[section]
\newtheorem{proposition}[theorem]{Proposition}
\newtheorem{lemma}[theorem]{Lemma}
\newtheorem{corollary}[theorem]{Corollary}
\theoremstyle{definition}
\newtheorem{definition}[theorem]{Definition}
\newtheorem{example}[theorem]{Example}
\theoremstyle{remark}
\newtheorem{remark}[theorem]{Remark}

\newcommand{\D}{\mathsf{D}}
\newcommand{\TV}{\mathrm{TV}}
\newcommand{\Reach}{\mathrm{Reach}}
\newcommand{\Lip}{\mathrm{Lip}}
\newcommand{\cP}{\mathcal{P}}
\newcommand{\cW}{\mathcal{W}}
\newcommand{\cS}{\mathcal{S}}
\newcommand{\cA}{\mathcal{A}}
\newcommand{\supp}{\operatorname{supp}}
\newcommand{\conv}{\operatorname{conv}}
\newcommand{\cone}{\operatorname{cone}}
\newcommand{\mass}{\operatorname{mass}}

\title{Precise Attribution of Model Interventions:\\
Attribution Sets, Mechanism Reach and Path Defects\\[4pt]
\large A companion to \emph{The Three-Level Posterior Tower}}
\author{Vinay K. Awasthi\\
Department of Applied Mathematics and Statistics\\
Johns Hopkins University\thanks{Verification scripts: \texttt{provenance.py},
\texttt{check\_attribution.py}, \texttt{provenance\_demo.py}; see \cref{sec:code}.}}
\date{October 4, 2026, draft 1}

\begin{document}
\maketitle

\begin{abstract}
A language model is trained on a corpus, extended with a second corpus, edited directly in its
weights, and tuned with reinforcement learning from human feedback. Each step changes some outputs
and leaves others alone, and new outputs appear. We ask three questions about such a history: what
explains each observed change, how precisely the data determine that explanation, and whether the
answer depends on the order of the steps. We answer them with three notions. An \emph{attribution
set} collects every decomposition of an observed output distribution into declared sources and an
unmodeled remainder; it is a polytope, and its emptiness, its range of unmodeled weight and its
certificates are computed by linear programming. \emph{Precision} is the width of the attribution
set; it is monotone along a ladder from exact sources to sources known by their supports to a
positive budget for unmodeled sources, and any slack in that budget forces imprecision.
\emph{Mechanism reach} declares what an intervention type may change: we prove that the reach set of
KL-regularized reinforcement learning with a bounded reward range is a polytope, so that attributing
a change to it is a linear program with an exact certificate, whereas a fixed reward with unknown
temperature gives a curve, and an unconstrained edit can explain anything. Finally, interventions are
compared quantitatively: comparisons compose by the triangle inequality and by a Lipschitz rule, with
exact Lipschitz constants for mixing and bounds for tilts, and a weight-level intervention may have no
Lipschitz constant on observed outputs at all. Order and loop defects measure path dependence, and
coarse evaluations certify fine differences. A synthetic four-stage experiment produces a full
provenance report. Every theorem is checked in exact arithmetic or with exactly verified linear
programming certificates.
\end{abstract}

\tableofcontents

%=====================================================================
\section{Introduction}\label{sec:intro}
%=====================================================================

Consider the life of a model. Corpus $C_1$ produces model $M_1$. Adding corpus $C_2$ produces $M_2$,
which changes some outputs, leaves others unchanged, and makes outcomes probable that were not before.
An engineer then edits weights directly, producing $M_3$, and reinforcement learning from human
feedback (RLHF) produces $M_4$. For such a history we want answers to three questions.

\begin{enumerate}[label=(Q\arabic*),leftmargin=*]
\item \emph{Attribution.} Which declared source or intervention explains each observed change, and
  what is left unexplained?
\item \emph{Precision.} How tightly do the observations determine that explanation?
\item \emph{Path dependence.} Does the result depend on the order of interventions, and does a
  sequence that should cancel actually cancel?
\end{enumerate}

\paragraph{Architecture.} The companion paper \cite{Awasthi-Tower} develops the \emph{local} structure:
an ambient probability monad $\D$, and over each outcome base $B$ a context comonad, a mixed
distributive law, admissibility data and an obstruction tower. Its \S5 establishes identifiability and
contamination results for a single output distribution. The present paper builds the
\emph{transport} half of the glue between local structures: interventions as maps between model
states, and quantitative comparisons between interventions. The other half of the glue, assembling
several contexts into a global state, is left to future work (\cref{sec:scope}).

\paragraph{Contributions.}
\begin{enumerate}[leftmargin=*]
\item Attribution sets for declared sources known exactly or by their supports, with an unmodeled
  budget; their structure, emptiness criterion and certificates (\cref{thm:attset}).
\item Precision as width, a monotone precision ladder separating structural, informational and
  unmodeled imprecision (\cref{thm:ladder}), and the fact that budget slack forces imprecision
  (\cref{prop:slack}).
\item Mechanism classes and a change-decomposition theorem with certificates (\cref{thm:change}); the
  reach set of KL-regularized RLHF with bounded rewards is a polytope (\cref{thm:tilt-polytope}),
  while a fixed reward with unknown temperature is not convex (\cref{prop:curve}).
\item Quantitative comparison of interventions with vertical and horizontal composition rules and an
  accumulation bound (\cref{thm:composition}); exact or bounded Lipschitz constants for mixing, tilts
  and coarsening, and the absence of any observable Lipschitz constant for some weight-level
  interventions (\cref{prop:lipschitz,prop:nonfaithful}); order and loop defects; monotonicity under
  coarsening (\cref{prop:coarse}).
\item A synthetic four-stage experiment with a complete provenance report (\cref{sec:demo}).
\end{enumerate}

\paragraph{What is not claimed.} We do not extract hypotheses from weights: declared sources must
come from reference models, checkpoints or ablations. All guarantees concern a finite probe set.
Attribution relative to a declared mechanism is not identification of a cause; causal statements
require running the compared paths. We treat output distributions as known; finite-sample precision
is an open problem (\cref{sec:scope}).

%=====================================================================
\section{Model states and observation}\label{sec:states}
%=====================================================================

\begin{definition}[Observable states]\label{def:obs}
Fix a finite set $X$ of \emph{probes} (prompts or evaluation contexts) and, for each $x\in X$, a
finite set $B_x$ of outcomes, which may be coarse classes of completions. The \emph{observable state
space} is $\cP=\prod_{x\in X}\D(B_x)$, with
\[
d(p,q)=\max_{x\in X}\TV(p_x,q_x),\qquad \TV(p_x,q_x)=\tfrac12\textstyle\sum_b\lvert p_x(b)-q_x(b)\rvert .
\]
\end{definition}

\begin{definition}[Model states and interventions]
A set $\cW$ of \emph{model states} (for instance weight vectors) carries an \emph{observation map}
$\Phi\colon\cW\to\cP$ sending a model to its output distributions on the probes. An
\emph{intervention} is a map $T\colon\cW\to\cW$; its \emph{observable effect} on $M$ is the pair
$(\Phi M,\Phi TM)$.
\end{definition}

\begin{remark}[Observation is not faithful]
Different models can have the same outputs on every probe, and the same intervention can then
produce different outputs. This is why statements about interventions on weights cannot in general be
reduced to statements about observable states (\cref{prop:nonfaithful}).
\end{remark}

\begin{definition}[New outcomes]
A softmax output assigns positive probability to every outcome, so ``a new outcome appears'' is made
precise by a threshold $\varepsilon>0$: outcome $b$ \emph{appears at level $\varepsilon$} at probe
$x$ under $M\to M'$ if $(\Phi M)_x(b)\le\varepsilon<(\Phi M')_x(b)$.
\end{definition}

From \cref{sec:attribution} to \cref{sec:mechanisms} we fix a probe, write $B$ for its finite outcome
set and $p\in\D(B)$ for an observed output distribution. All statements apply probe by probe.

%=====================================================================
\section{Attribution sets and precision}\label{sec:attribution}
%=====================================================================

\begin{definition}[Declared sources and attribution sets]\label{def:attset}
A \emph{declared source} is either \emph{exact}, a distribution $A_i\in\D(B)$, or \emph{known by its
support}, a set $S_i\subseteq B$. An \emph{unmodeled budget} is a number $\bar\mu\in[0,1]$. A
\emph{decomposition} of $p$ is a family of nonnegative measures $s_1,\dots,s_k$ on $B$, the
\emph{contributions}, such that
\begin{enumerate}[label=(\alph*)]
\item $s_i=\lambda_iA_i$ for some $\lambda_i\ge0$ if source $i$ is exact, and $\supp s_i\subseteq S_i$
  if source $i$ is known by its support, in which case $\lambda_i=\mass(s_i)$;
\item $\sum_is_i\le p$ pointwise, and $\mu:=1-\sum_i\lambda_i\le\bar\mu$.
\end{enumerate}
The \emph{unmodeled part} is $p-\sum_is_i=\mu U$ with $U\in\D(B)$ when $\mu>0$. The \emph{attribution
set} $W_{\bar\mu}(p)$ is the set of all decompositions. The \emph{minimal unmodeled weight}
$\mu_{\min}(p)$ is the least $\mu$ over decompositions with $\bar\mu=1$.
\end{definition}

The model is Huber's contamination model \cite{Huber1964} with declared sources. A decomposition
determines $U$, so the attribution set records exactly the explanations compatible with $p$.

\begin{theorem}[Structure of attribution sets]\label{thm:attset}
\leavevmode
\begin{enumerate}[label=(\roman*)]
\item $W_{\bar\mu}(p)$ is a polytope, cut out by linear constraints on the contributions.
\item $W_{\bar\mu}(p)\ne\emptyset$ if and only if $\mu_{\min}(p)\le\bar\mu$.
\item If $W_{\bar\mu}(p)\ne\emptyset$, the unmodeled weights it contains form the interval
  $[\mu_{\min}(p),\bar\mu]$.
\item (Certificate) Suppose $y\in\mathbb{R}^B_{\ge0}$ satisfies $\langle y,a\rangle\ge1$ for every
  admissible version $a$ of every source, that is $\langle y,A_i\rangle\ge1$ for exact sources and
  $y_b\ge1$ for $b\in S_i$ for support-known sources. Then every decomposition has
  $\mu\ge1-\langle y,p\rangle$, and the best such bound equals $\mu_{\min}(p)$.
\end{enumerate}
\end{theorem}
\begin{proof}
(i) Constraints (a) and (b) are linear in the variables $\lambda_i$ (exact sources) and $s_i$
(support-known sources), with $s_i\ge0$, and the feasible set is bounded since $\sum_is_i\le p$.
(ii) and (iii) If a decomposition has weight vector $\lambda$, then so does $t\lambda$ (with scaled
contributions) for every $t\in[0,1]$, which raises $\mu$ continuously up to $1$. Hence the attainable
$\mu$ form $[\mu_{\min}(p),1]$, and the budget cuts this to $[\mu_{\min}(p),\bar\mu]$.
(iv) Each contribution is a nonnegative multiple $\lambda_ia$ of an admissible version $a$ of its
source, so $\langle y,\sum_is_i\rangle\ge\sum_i\lambda_i=1-\mu$, while $\langle y,\sum_is_i\rangle
\le\langle y,p\rangle$. Optimality of the best bound is linear programming duality \cite{Schrijver1986}:
$1-\mu_{\min}(p)$ is the value of $\max\{\sum_i\lambda_i\}$ over (a), (b), and its dual is
$\min\{\langle y,p\rangle\}$ under the stated constraints on $y$.
\end{proof}

\begin{definition}[Precision]\label{def:precision}
For a nonempty attribution set, the \emph{width} of source $i$ is $w_i=\max\lambda_i-\min\lambda_i$
over $W_{\bar\mu}(p)$, and the width of the unmodeled weight is $w_\mu=\max\mu-\min\mu$. For an outcome
$b$ with $p(b)>0$, the \emph{observation-level attribution} $\pi_i(b)=s_i(b)/p(b)$ is the probability
that an observation of $b$ came from source $i$; its width is defined likewise. The attribution is
\emph{exact} if all widths vanish. All widths are values of linear programs over $W_{\bar\mu}(p)$.
\end{definition}

\begin{theorem}[Monotone precision ladder]\label{thm:ladder}
Enlarging the admissible versions of any source, or raising the budget, can only enlarge the
attribution set; hence all widths are nondecreasing. In particular, writing $\mu^{\mathrm{ex}}_{\min}$
for the minimal unmodeled weight with exact sources, the attribution sets
\[
W^{(1)}=W^{\mathrm{exact}}_{\mu^{\mathrm{ex}}_{\min}}(p)\ \subseteq\
W^{(2)}=W^{\mathrm{support}}_{\mu^{\mathrm{ex}}_{\min}}(p)\ \subseteq\
W^{(3)}=W^{\mathrm{support}}_{\bar\mu}(p)\qquad(\bar\mu\ge\mu^{\mathrm{ex}}_{\min})
\]
are nonempty and nested, and each width satisfies $w^{(1)}\le w^{(2)}\le w^{(3)}$.
\end{theorem}
\begin{proof}
Every decomposition with exact sources is one with the corresponding supports, since $\supp(\lambda_i
A_i)\subseteq\supp A_i$; and every decomposition within budget $\bar\mu$ is one within any larger
budget. The width of a coordinate over a larger set is at least its width over a smaller set.
$W^{(1)}$ is nonempty by \cref{thm:attset}(ii).
\end{proof}

We call the increments $w^{(1)}$, $w^{(2)}-w^{(1)}$ and $w^{(3)}-w^{(2)}$ the \emph{structural},
\emph{informational} and \emph{unmodeled} imprecision. They are nonnegative by \cref{thm:ladder}; we do
not claim that they combine additively across sources.

\begin{proposition}[Budget slack forces imprecision]\label{prop:slack}
If $\mu_{\min}(p)<\bar\mu$ and $\mu_{\min}(p)<1$, then $w_\mu=\bar\mu-\mu_{\min}(p)>0$, and every source
whose maximal weight is positive has positive width.
\end{proposition}
\begin{proof}
The first claim is \cref{thm:attset}(iii). Take a decomposition $\lambda$ with $\mu=\mu_{\min}(p)$.
For $t=(1-\bar\mu)/(1-\mu_{\min}(p))<1$, the scaled family $t\lambda$ is again a decomposition, since its
unmodeled weight is $\bar\mu$. If $\lambda_i>0$, the two decompositions give different values of the
$i$-th weight. If $\lambda_i=0$, this value differs from the positive maximal weight. Either way
$w_i>0$.
\end{proof}

So precision and robustness trade off: a budget for unmodeled sources protects against rejecting the
declared model, and every unit of slack in that budget is paid for in width.

\begin{proposition}[When attribution is exact]\label{prop:exact}
With exact sources and budget $\mu_{\min}(p)=0$, the attribution is exact for every $p$ in
$\conv(A_1,\dots,A_k)$ if and only if the sources are affinely independent. For any sources, if the
exclusive region $R_i=\{b:A_i(b)>0,\ A_j(b)=0\ (j\ne i)\}$ has $A_i(R_i)>0$, then
$\lambda_i=p(R_i)/A_i(R_i)$ in every decomposition, and $\pi_i(b)=\lambda_iA_i(b)/p(b)$.
\end{proposition}
\begin{proof}
This is \cite[Prop.~5.7]{Awasthi-Tower}.
\end{proof}

%=====================================================================
\section{Mechanism classes and attribution of change}\label{sec:mechanisms}
%=====================================================================

\begin{definition}[Mechanism class]
A \emph{mechanism class} for an intervention type assigns to each $p\in\D(B)$ a set
$\Reach(p)\subseteq\D(B)$ of output distributions the intervention may produce from $p$. A change
$p\to p'$ is \emph{reachable} if $p'\in\Reach(p)$.
\end{definition}

\begin{example}[Three polyhedral classes]\label{ex:classes}
\leavevmode
\begin{enumerate}[label=(\alph*)]
\item \emph{Source expansion} (adding a corpus): $\Reach(p)=\conv(\cA\cup\mathcal{N})$, where $\cA$
  are the old sources and $\mathcal{N}$ the new ones.
\item \emph{Locality claim} for an edit with declared scope $X_0\subseteq X$ and tolerance $\tau$: at
  a probe outside $X_0$, $\Reach(p)=\{q:\TV(q,p)\le\tau\}$.
\item \emph{Unconstrained edit}: $\Reach(p)=\D(B)$.
\end{enumerate}
\end{example}

\begin{proposition}[KL-regularized reinforcement learning is an exponential tilt]\label{prop:gibbs}
For a reward $r\colon B\to\mathbb{R}$ and $\beta>0$, the maximizer of
$\mathbb{E}_q[r]-\beta\,\mathrm{KL}(q\,\|\,p)$ over $q\in\D(B)$ is
\[
T_r(p)(b)=\frac{p(b)\,e^{r(b)/\beta}}{\sum_cp(c)\,e^{r(c)/\beta}} .
\]
\end{proposition}
\begin{proof}
For $q\ll p$, $\mathbb{E}_q[r]-\beta\,\mathrm{KL}(q\|p)=\beta\log Z-\beta\,\mathrm{KL}(q\|T_r(p))$ with
$Z=\sum_cp(c)e^{r(c)/\beta}$, which is maximal exactly at $q=T_r(p)$; and $q\not\ll p$ gives
$-\infty$. This is the standard Gibbs variational principle, used for RLHF in
\cite{Korbak2022,Rafailov2023}.
\end{proof}

Writing $w(b)=e^{(r(b)-r_{\min})/\beta}$, a reward with values in $[r_{\min},r_{\max}]$ gives weights
in $[1,K]$ with $K=e^{(r_{\max}-r_{\min})/\beta}$, and $T_r(p)$ depends only on $w$. The tilt is a
reweighting by a likelihood, the operation called conditioning in \cite[\S4]{Awasthi-Tower}; it is not
conditioning on an observed event. It preserves support exactly, $\supp T_r(p)=\supp p$, but it can
raise a small probability by a factor up to $K$.

\begin{theorem}[The tilt reach set is a polytope]\label{thm:tilt-polytope}
Let $\Reach_K(p)=\{T_r(p):r(B)\subseteq[r_{\min},r_{\max}]\}$ with $K$ as above. Then
\[
\Reach_K(p)=\big\{q\in\D(B):\ \supp q=\supp p,\ \ q_b\,p_c\le K\,q_c\,p_b\ \text{ for all }b,c\in\supp p\big\},
\]
a convex polytope. Equivalently $q\in\Reach_K(p)$ iff $\supp q=\supp p$ and
\[
K^*(p,q)=\frac{\max_{b\in\supp p}q_b/p_b}{\min_{b\in\supp p}q_b/p_b}\le K,
\]
so $\beta\log K^*(p,q)$ is the least reward range that explains the change.
\end{theorem}
\begin{proof}
If $q=T_r(p)$, then $q_b/p_b=w_b/Z$ for $b\in\supp p$, so the ratios lie within a factor $K$ of each
other, which is the stated inequality. Conversely, if $q$ satisfies the inequalities and has support
$\supp p$, put $w_b=(q_b/p_b)/\min_c(q_c/p_c)\in[1,K]$ for $b\in\supp p$ and $w_b=1$ elsewhere; then
$T_w(p)=q$. The inequalities are linear in $q$ for fixed $p$; together with $\sum_bq_b=1$ and $q_b=0$
off $\supp p$ they define a polytope. A point satisfying them cannot vanish anywhere on $\supp p$:
if $q_b=0$ then $q_c\,p_b\le K\,q_b\,p_c=0$ for all $c$, contradicting $\sum q=1$. So the constraint
$\supp q=\supp p$ is implied, and the set is closed.
\end{proof}

\begin{proposition}[Fixed reward, unknown temperature]\label{prop:curve}
For a fixed reward $r$, the set $\{T_{r/\beta'}(p):\beta'>0\}$ is a curve, and in general not convex:
for $p$ uniform on three outcomes and $r=(0,1,2)$, the midpoint of the points with $1/\beta'=1$ and
$1/\beta'=3$ is not on the curve.
\end{proposition}
\begin{proof}
On the curve, $\log(q_b/p_b)-\log(q_0/p_0)$ is proportional to $r_b-r_0$, so $\log(q_2/q_0)=
2\log(q_1/q_0)$ for this $p$ and $r$; the midpoint violates this identity, as checked in
\cref{sec:code}.
\end{proof}

So the geometry of a reach set depends on what is assumed about the mechanism: a bounded reward range
gives a polytope, a fixed reward gives a curve, and no assumption gives everything.

\begin{definition}[Residual of a change]
Let $\cone\Reach(p)=\{tq:t\ge0,\ q\in\Reach(p)\}$. An \emph{explained part} of a change $p\to p'$ is a
measure $s\in\cone\Reach(p)$ with $s\le p'$; the remainder $p'-s$ is the \emph{residual}. The
\emph{residual weight} of the change is
\[
\mu^T_{\min}(p,p')=1-\max\{\mass(s):\ s\in\cone\Reach(p),\ s\le p'\}.
\]
\end{definition}

\begin{theorem}[Change decomposition]\label{thm:change}
Let $\Reach(p)$ be closed and convex.
\begin{enumerate}[label=(\roman*)]
\item The maximum defining $\mu^T_{\min}(p,p')$ is attained, and $\mu^T_{\min}(p,p')=0$ if and only if
  $p'\in\Reach(p)$.
\item (Certificate) If $y\ge0$ satisfies $\langle y,q\rangle\ge1$ for all $q\in\Reach(p)$, then
  $\mu^T_{\min}(p,p')\ge1-\langle y,p'\rangle$. If $\cone\Reach(p)=\{s\ge0:Gs\le0\}$ is polyhedral, the
  best such bound is attained by a pair $y\ge0$, $\nu\ge0$ with $y+G^{\top}\nu\ge\mathbf{1}$, and it equals
  $\mu^T_{\min}(p,p')$.
\item (Precision of change attribution) The optimal explained parts form a polytope when the cone is
  polyhedral, and the range of each $s(b)$ over it is given by linear programs.
\end{enumerate}
\end{theorem}
\begin{proof}
(i) $\Reach(p)\subseteq\D(B)$ is compact and avoids $0$, so its cone is closed; intersected with
$\{0\le s\le p'\}$ it is compact, and $\mass$ is continuous. If the maximum is $1$, then $s\le p'$ and
$\mass(s)=\mass(p')$ force $s=p'$, so $p'\in\cone\Reach(p)$ with mass one, that is, $p'\in\Reach(p)$.
(ii) For $s=tq$ with $q\in\Reach(p)$, $\mass(s)=t\le t\langle y,q\rangle=\langle y,s\rangle\le\langle
y,p'\rangle$. In the polyhedral case, $\max\{\mathbf 1^{\top}s: Gs\le0,\ s\le p',\ s\ge0\}$ has dual
$\min\{\langle y,p'\rangle:y+G^{\top}\nu\ge\mathbf 1,\ y,\nu\ge0\}$, and strong duality holds
\cite{Schrijver1986}; every dual-feasible pair satisfies $\langle y,s\rangle\ge\mathbf 1^{\top}s-\langle
\nu,Gs\rangle\ge\mathbf 1^{\top}s$ on the cone. (iii) The optimal set is a face of a polytope.
\end{proof}

\begin{remark}[What residuals mean]
A positive residual weight says that at least this much of the observed change is not produced by any
intervention of the declared type. It does not say what produced it. For the tilt class the cone is
$\{s\ge0:\ s_b\,p_c\le K\,s_c\,p_b,\ s=0\text{ off }\supp p\}$, so attributing a change to RLHF is a
linear program with an exact certificate. For an unconstrained edit the residual is always zero: an
edit can be tested only through a claim about it, such as locality.
\end{remark}

%=====================================================================
\section{Comparing interventions: composition and path defects}\label{sec:composition}
%=====================================================================

\begin{definition}[Quantitative comparison]\label{def:comparison}
For interventions $F,G$ and a set $\cS\subseteq\cW$ of model states,
\[
d_\cS(F,G)=\sup_{M\in\cS}d(\Phi FM,\Phi GM).
\]
This is a pseudometric on interventions. We write $F\approx_\varepsilon G$ on $\cS$ if
$d_\cS(F,G)\le\varepsilon$, and call this a \emph{quantitative comparison}.
\end{definition}

A quantitative comparison records how far two interventions are apart; it is not by itself a 2-cell
in a 2-category, and we do not call it one. \Cref{cor:enriched} identifies the structure in which it is
the hom-metric of an enriched category.

\begin{definition}[Lipschitz on observables]
An intervention $G$ is \emph{$L$-Lipschitz on observables} over $\mathcal{T}\subseteq\cW$ if
$d(\Phi GM,\Phi GM')\le L\,d(\Phi M,\Phi M')$ for all $M,M'\in\mathcal T$. An \emph{output-level
mechanism} is a map $T\colon\cP\to\cP$ acting on observable states; for it the Lipschitz constant is the
usual one for $d$.
\end{definition}

\begin{proposition}[Lipschitz constants of output-level mechanisms]\label{prop:lipschitz}
\leavevmode
\begin{enumerate}[label=(\roman*)]
\item Mixing with a fixed target, $p\mapsto(1-a)p+a\,t$, is exactly $(1-a)$-Lipschitz.
\item A tilt with weights in $[1,K]$ is $\tfrac{3K-1}{2}$-Lipschitz, and no constant below $K$ works in
  general.
\item Coarsening along any map of outcomes is $1$-Lipschitz.
\end{enumerate}
\end{proposition}
\begin{proof}
(i) The difference of images is $(1-a)(p-q)$. (iii) Data processing for total variation. (ii) Let
$Z=\langle w,p\rangle\ge1$ and $Z'=\langle w,q\rangle$. Then
$T(p)-T(q)=w\odot(p-q)/Z+w\odot q\,(1/Z-1/Z')$, so
$\lVert T(p)-T(q)\rVert_1\le\big(K\lVert p-q\rVert_1+\lvert Z-Z'\rvert\big)/Z\le K\lVert
p-q\rVert_1+\lvert Z-Z'\rvert$, using $\lVert w\odot q\rVert_1=Z'$ and $Z\ge1$. Since $p-q$ has total mass zero,
$\lvert Z-Z'\rvert=\lvert\langle w-\tfrac{K+1}2\mathbf 1,p-q\rangle\rvert\le\tfrac{K-1}2\lVert
p-q\rVert_1$. For the lower bound take $p=(\epsilon,1-\epsilon)$ and $w=(K,1)$: the first coordinate
of $T(p)$ is $K\epsilon/(1+(K-1)\epsilon)$, whose derivative at $\epsilon=0$ is $K$.
\end{proof}

So tilts can amplify differences, which matters for composition.

\begin{proposition}[No observable Lipschitz constant]\label{prop:nonfaithful}
If $\Phi M=\Phi M'$ but $\Phi GM\ne\Phi GM'$, then $G$ is not $L$-Lipschitz on observables over
$\{M,M'\}$ for any $L$.
\end{proposition}
\begin{proof}
The left side of the Lipschitz inequality is positive and the right side is zero.
\end{proof}

A weight edit that depends on parts of the model not visible on the probes is of this kind. For such
interventions the composition bound below must use a metric on $\cW$, and conclusions at the probe
level are conditional on that metric.

\begin{theorem}[Composition of quantitative comparisons]\label{thm:composition}
\leavevmode
\begin{enumerate}[label=(\roman*)]
\item (Vertical) $d_\cS(F,H)\le d_\cS(F,G)+d_\cS(G,H)$.
\item (Horizontal) If $G'$ is $L$-Lipschitz on observables over $F(\cS)\cup F'(\cS)$, then
  \[
  d_\cS(GF,G'F')\le d_{F(\cS)}(G,G')+L\,d_\cS(F,F').
  \]
\item (Accumulation) For paths $F_n\cdots F_1$ and $G_n\cdots G_1$, let $\cS_0=\cS$ and
  $\cS_j=F_j(\cS_{j-1})$, put $\varepsilon_j=d_{\cS_{j-1}}(F_j,G_j)$, and suppose each $G_j$ is
  $L_j$-Lipschitz on observables over the states reached by both paths. Then
  \[
  \delta_j=d_\cS(F_j\cdots F_1,\,G_j\cdots G_1)\quad\text{satisfies}\quad
  \delta_j\le\varepsilon_j+L_j\delta_{j-1},
  \]
  and hence
  \[
  \delta_n\le\sum_{j=1}^n\varepsilon_j\prod_{i=j+1}^nL_i .
  \]
\end{enumerate}
\end{theorem}
\begin{proof}
(i) is the triangle inequality for $d$. (ii) For $M\in\cS$,
\[
d(\Phi GFM,\Phi G'F'M)\le d(\Phi GFM,\Phi G'FM)+d(\Phi G'FM,\Phi G'F'M);
\]
the first term is at most $d_{F(\cS)}(G,G')$ and the second at most $L\,d(\Phi FM,\Phi F'M)$. (iii) Apply (ii) with $F=F_j$, $G=G_j$ to the states $F_{j-1}\cdots F_1M$
and $G_{j-1}\cdots G_1M$, and unroll the recursion.
\end{proof}

\begin{corollary}[Enriched structure]\label{cor:enriched}
Output-level mechanisms on $\cP$ with Lipschitz constant at most $1$, with $d=d_{\cP}$, form a
category enriched in pseudometric spaces in Lawvere's sense \cite{Lawvere1973}:
$d(\mathrm{id},\mathrm{id})=0$ and $d(GF,G'F')\le d(G,G')+d(F,F')$. Mixings and coarsenings belong to it; tilts do not, since their
Lipschitz constant can exceed $1$, so in general the structure is the Lipschitz-weighted one of
\cref{thm:composition}.
\end{corollary}

\begin{definition}[Order and loop defects]
For interventions $E$ and $R$, the \emph{order defect} on $\cS$ is $d_\cS(RE,ER)$. For a sequence
of interventions $T_1,\dots,T_n$, the \emph{loop defect} is $d_\cS(T_n\cdots T_1,\,\mathrm{id})$.
\end{definition}

A positive order defect certifies that the two orders produce different observable models on the
probes; it becomes a causal statement once both orders have actually been run. We call the loop defect
a defect rather than a holonomy, since no connection has been defined.

\begin{proposition}[Coarse evaluation]\label{prop:coarse}
For any coarsening $g$ of outcomes, $\TV(g_*p,g_*q)\le\TV(p,q)$. Hence every order, loop or comparison
defect measured on a coarse evaluation is a lower bound for the fine one, and a positive coarse defect
certifies a positive fine defect.
\end{proposition}

%=====================================================================
\section{A four-stage provenance experiment}\label{sec:demo}
%=====================================================================

We run a synthetic history through all four stages and report what the framework certifies. The
experiment is in \texttt{provenance\_demo.py}; all distributions are exact rationals.

\paragraph{Setup.} There are three probes $x_1,x_2,x_3$ and six outcomes $y_1,\dots,y_6$. Corpus $C_1$
has two sources $A_1,A_2$, and corpus $C_2$ has one source $N$, which is the only one giving mass to
$y_6$:
\begin{center}\small
\begin{tabular}{@{}lccc@{}}
\toprule
probe & $A_1$ & $A_2$ & $N$\\
\midrule
$x_1$ & $(.4,.3,.2,.1,0,0)$ & $(0,.1,.2,.3,.4,0)$ & $(0,0,0,.2,.3,.5)$\\
$x_2$ & $(.1,.1,.4,.4,0,0)$ & $(.3,.3,.1,.1,.2,0)$ & $(0,0,0,0,.5,.5)$\\
$x_3$ & $(.5,.5,0,0,0,0)$ & $(0,0,.5,.5,0,0)$ & $(0,0,0,0,.5,.5)$\\
\bottomrule
\end{tabular}
\end{center}
$M_1$ mixes $A_1,A_2$ with weights $(0.6,0.4)$ at every probe. $M_2$ uses weights $(0.45,0.25,0.30)$ at
$x_1$, $(0.5,0.3,0.2)$ at $x_2$, and leaves $x_3$ unchanged. The edit producing $M_3$ declares scope
$\{x_1\}$ and moves $20\%$ of the mass at $x_1$ to $y_3$, but also leaks $5\%$ to $y_1$ at $x_2$. RLHF
producing $M_4$ is a tilt with weights $(1,2,4,1,3,1)$, so $K=4$; a \emph{drifted} variant additionally
multiplies $y_6$ by $10$ at $x_2$.

\paragraph{Stage 2: adding $C_2$.} With exact sources the attribution is exact at every probe: the
weights $(A_1,A_2,N)$ are recovered as $(0.45,0.25,0.30)$, $(0.5,0.3,0.2)$ and $(0.6,0.4,0)$, with widths
zero. Outcome $y_6$ appears at level $\varepsilon=0.01$ at $x_1$ and $x_2$ and not at $x_3$, where
$\TV(M_1,M_2)=0$. Since $y_6$ is exclusive to $N$, \cref{prop:exact} gives $\lambda_N=p(y_6)/N(y_6)$,
namely $3/10$ and $1/5$. An observation of $y_5$ at $x_1$ came from $A_2$ with probability $10/19$ and
from $N$ with probability $9/19$. The precision ladder at $x_1$ is:
\begin{center}\small
\begin{tabular}{@{}lcccc@{}}
\toprule
stage & $A_1$ & $A_2$ & $N$ & $\mu$\\
\midrule
exact sources, budget $0$ & $0.45$ & $0.25$ & $0.30$ & $0$\\
supports only, budget $0$ & $[0.18,0.66]$ & $[0,0.67]$ & $[0.15,0.52]$ & $0$\\
supports only, budget $0.1$ & $[0.08,0.66]$ & $[0,0.67]$ & $[0.05,0.52]$ & $[0,0.1]$\\
\bottomrule
\end{tabular}
\end{center}
The structural imprecision is zero, the informational imprecision is large, and the budget lowers the
lower ends further, as \cref{thm:ladder,prop:slack} predict. The lower bound $0.18$ for $A_1$ is the mass of
its exclusive outcome $y_1$.

\paragraph{Stage 3: the edit.} The locality test with $\tau=0.01$ passes at $x_3$ ($\TV=0$) and fails at
$x_2$ ($\TV(M_2,M_3)=0.043$), detecting the leak. Relative to the corpus sources $A_1,A_2,N$, the
certified unmodeled weight is $0.20$ at $x_1$ and $0.05$ at $x_2$, exactly the mass moved by the edit
and by the leak, with certificates $y$ satisfying $\min_i\langle y,A_i\rangle=1$ and $\langle
y,p\rangle=0.80$, respectively $0.95$. At $x_1$ the total-variation change is $0.172$, smaller than the
certified unexplained weight $0.20$: the two quantities answer different questions.

\paragraph{Stage 4: RLHF.} For the pure tilt, the implied ratio is $K^*=4$ at every probe, so every
change is reachable with the claimed reward range. For the drifted variant, $K^*=10$ at $x_2$, so the
change is not reachable; the residual weight relative to the tilt class is $0.1934$, with an exactly
verified certificate. The other probes remain reachable, so the departure from the declared mechanism
is localized.

\paragraph{Paths.} At $x_1$, with $E$ the edit and $R$ the tilt, the order defect is
$\TV(REp,ERp)=51/518\approx0.0985$, so the order matters. Tilting by $w$ and then by $1/w$ has loop defect
zero. Comparing $E$ with a weaker edit $E'$ (moving $25\%$ instead of $20\%$),
$\TV(REp,RE'p)=0.0479$ is within the bound $\tfrac{3K-1}2\TV(Ep,E'p)=5.5\times0.043$, with observed
amplification $1.11$. On the coarse evaluation $\{y_1,y_2\},\{y_3,y_4\},\{y_5,y_6\}$ the order defect is
$0.0858$, smaller but still positive, so the coarse evaluation alone certifies that the order matters.

\paragraph{Summary.} The report separates what changed, what the declared sources and mechanisms
explain, how precisely, and where they fail:
\begin{center}\small
\begin{tabular}{@{}llll@{}}
\toprule
step & behavioral change & attribution & certified residual\\
\midrule
add $C_2$ & $\TV=0.18$ at $x_1,x_2$; $0$ at $x_3$ & exact; $\lambda_N=0.3,\,0.2,\,0$ & none\\
edit & $\TV=0.172$ at $x_1$; leak $0.043$ at $x_2$ & locality fails at $x_2$ & $0.20$ ($x_1$), $0.05$ ($x_2$)\\
RLHF & tilt, $K^*=4$ & reachable at all probes & none\\
RLHF, drifted & $K^*=10$ at $x_2$ & not reachable at $x_2$ & $0.19$ ($x_2$)\\
\bottomrule
\end{tabular}
\end{center}

%=====================================================================
\section{Relation to the posterior tower and scope}\label{sec:scope}
%=====================================================================

\paragraph{Local structure and glue.} In the architecture of \cite{Awasthi-Tower}, the probability monad
$\D$ is ambient, and each base $B$ carries local structure: a context comonad, a distributive law,
admissibility data and an obstruction tower. This paper supplies the transport between local
structures: interventions, mechanism classes and quantitative comparison. Deterministic base change is
already a structure-preserving map in \cite[Prop.~9.13]{Awasthi-Tower}; extending it to Markov kernels
and to interventions observed through $\Phi$ is the next step.

\paragraph{Structural provenance.} The obstruction tower of \cite[\S10]{Awasthi-Tower} gives a second
coordinate. For a claim $c$ attached to a model, let $k(c)$ be its first failing level. Along an
intervention, the pair $(k_{\mathrm{before}},k_{\mathrm{after}})$ records which level of admissibility
changed. Because failure is monotone up the tower \cite[Thm.~10.5]{Awasthi-Tower}, a transition
$k:1\to2$ means that a level-one obstruction was removed while a level-two one remains; it does not
mean that a level-two obstruction was created. We leave the systematic study of such transitions to
future work.

\paragraph{Open problems.}
\begin{enumerate}[leftmargin=*]
\item \emph{Finite-sample precision.} Output distributions are estimated from samples. Replacing $p$ by
  a confidence region $\mathcal C_\delta$ turns each linear program into a robust program and each width
  into a statistical width.
\item \emph{Weight-level metrics.} Choose a metric on $\cW$ for which real interventions are Lipschitz,
  so that \cref{thm:composition} applies beyond output-level mechanisms.
\item \emph{Higher levels.} Lift comparisons to second- and third-order states with the Kantorovich
  metric \cite{FritzPerrone2019}, and relate them to the obstruction tower.
\item \emph{Multi-context assembly.} Glue several probes or contexts into one global state, with
  approximate compatibility measured by quantitative comparisons.
\item \emph{Empirical validation.} Run the provenance report on real checkpoints, with declared sources
  from reference models.
\end{enumerate}

\paragraph{Related work.} Attribution of predictions to training data is studied with influence
functions \cite{KohLiang2017}; locality claims for weight edits appear in \cite{Meng2022}; the view of
KL-regularized RLHF as Bayesian reweighting is in \cite{Korbak2022,Rafailov2023}; contamination models
go back to \cite{Huber1964}, and lower and upper probabilities to \cite{Dempster1967,Shafer1976}.
Sheaf-theoretic contextuality \cite{AbramskyBrandenburger2011} and its comonadic treatment of
simulations \cite{Abramsky2019} concern compatibility of local models and maps between them; to our
knowledge they provide neither quantitative comparisons between interventions nor the probability
hierarchy of \cite{Awasthi-Tower}.

\appendix

%=====================================================================
\section{Code and computational checks}\label{sec:code}
%=====================================================================

The library \texttt{provenance.py} implements attribution sets, widths, residuals and certificates.
The script \texttt{check\allowbreak\_attribution.py} verifies the following.
\cref{thm:attset} (nonemptiness criterion, the interval $[\mu_{\min},\bar\mu]$, and exactly verified dual
certificates); \cref{thm:ladder} (monotone widths along the ladder and in the budget);
\cref{prop:slack}; \cref{thm:tilt-polytope} in exact arithmetic (every tilt satisfies the inequalities,
every point satisfying them is reconstructed exactly as a tilt, and mixtures of tilts are tilts);
\cref{prop:curve}; \cref{thm:change} for the tilt class (exactly verified certificates $(y,\nu)$, and
zero residual exactly for reachable changes); \cref{prop:lipschitz} (the exact constant for mixing, the
tilt bound, and an example with amplification $4.0$ for $K=4$); \cref{prop:nonfaithful};
\cref{thm:composition}(i)--(iii) in exact arithmetic; and \cref{prop:coarse}. All 18 checks pass.
The script \texttt{provenance\allowbreak\_demo.py} produces the report of \cref{sec:demo}. Linear programs are solved in
floating point, and their certificates are rounded to rationals and verified exactly.

%=====================================================================
\begin{thebibliography}{99}

\bibitem{Abramsky2019}
S.~Abramsky, R.~S. Barbosa, M.~Karvonen, and S.~Mansfield.
\newblock A comonadic view of simulation and quantum resources.
\newblock In \emph{Proceedings of the 34th Annual ACM/IEEE Symposium on Logic in Computer Science
  (LICS)}, 2019.

\bibitem{AbramskyBrandenburger2011}
S.~Abramsky and A.~Brandenburger.
\newblock The sheaf-theoretic structure of non-locality and contextuality.
\newblock \emph{New Journal of Physics}, 13:113036, 2011.

\bibitem{Awasthi-Tower}
V.~K. Awasthi.
\newblock The three-level posterior tower: hierarchical Bayesian inference and mixed distributive
  laws.
\newblock Revised version 8, Johns Hopkins University, October 2026.

\bibitem{Dempster1967}
A.~P. Dempster.
\newblock Upper and lower probabilities induced by a multivalued mapping.
\newblock \emph{Annals of Mathematical Statistics}, 38(2):325--339, 1967.

\bibitem{FritzPerrone2019}
T.~Fritz and P.~Perrone.
\newblock A probability monad as the colimit of spaces of finite samples.
\newblock \emph{Theory and Applications of Categories}, 34(7):170--220, 2019.

\bibitem{Huber1964}
P.~J. Huber.
\newblock Robust estimation of a location parameter.
\newblock \emph{Annals of Mathematical Statistics}, 35(1):73--101, 1964.

\bibitem{KohLiang2017}
P.~W. Koh and P.~Liang.
\newblock Understanding black-box predictions via influence functions.
\newblock In \emph{Proceedings of the 34th International Conference on Machine Learning (ICML)}, 2017.

\bibitem{Korbak2022}
T.~Korbak, E.~Perez, and C.~L. Buckley.
\newblock {RL} with {KL} penalties is better viewed as {B}ayesian inference.
\newblock In \emph{Findings of the Association for Computational Linguistics: EMNLP}, 2022.

\bibitem{Lawvere1973}
F.~W. Lawvere.
\newblock Metric spaces, generalized logic, and closed categories.
\newblock \emph{Rendiconti del Seminario Matematico e Fisico di Milano}, 43:135--166, 1973.

\bibitem{Meng2022}
K.~Meng, D.~Bau, A.~Andonian, and Y.~Belinkov.
\newblock Locating and editing factual associations in {GPT}.
\newblock In \emph{Advances in Neural Information Processing Systems (NeurIPS)}, 2022.

\bibitem{Rafailov2023}
R.~Rafailov, A.~Sharma, E.~Mitchell, S.~Ermon, C.~D. Manning, and C.~Finn.
\newblock Direct preference optimization: your language model is secretly a reward model.
\newblock In \emph{Advances in Neural Information Processing Systems (NeurIPS)}, 2023.

\bibitem{Schrijver1986}
A.~Schrijver.
\newblock \emph{Theory of Linear and Integer Programming}.
\newblock Wiley, Chichester, 1986.

\bibitem{Shafer1976}
G.~Shafer.
\newblock \emph{A Mathematical Theory of Evidence}.
\newblock Princeton University Press, 1976.

\end{thebibliography}

\end{document}
